% Generated draft: requires mathematical and source review.
\documentclass[11pt]{article}
\usepackage[margin=25mm]{geometry}
\usepackage{fontspec}
\setmainfont{Latin Modern Roman}
\usepackage{amsmath,amssymb,mathtools,unicode-math}
\setmathfont{Latin Modern Math}
\usepackage{enumitem,needspace,xurl,hyperref}
\setlength{\parindent}{0pt}
\setlength{\parskip}{5pt}
\setlength{\emergencystretch}{3em}
\newcommand{\N}{\mathbb{N}}
\newcommand{\Z}{\mathbb{Z}}
\newcommand{\Q}{\mathbb{Q}}
\newcommand{\R}{\mathbb{R}}
\newcommand{\C}{\mathbb{C}}
\newcommand{\F}{\mathbb{F}}
\newcommand{\A}{\mathbb{A}}
\newcommand{\PP}{\mathbb{P}}
\newcommand{\E}{\mathbb{E}}
\newcommand{\eps}{\varepsilon}
\newcommand{\dd}{\,\mathrm{d}}
\begin{document}
\section*{UnsolvedMath ID: 800255}
\textbf{Source problem code:} catalogue entry
\subsection{Definitions and mathematical statement}
Let \(P=(P_1,P_2):\C^2\to\C^2\) be a polynomial map, with \(P_1,P_2\in\C[z_1,z_2]\). Its complex Jacobian determinant is
\[
J_P(z)=\det\left(\frac{\partial P_i}{\partial z_j}(z)\right)_{i,j=1}^{2}.
\]
Here, nonvanishing Jacobian means \(J_P(z)\ne0\) for every \(z\in\C^2\). A map is biholomorphic if it is bijective and both it and its inverse are holomorphic.

\textbf{Statement (supplied Problem 54).} For every polynomial map \(P:\C^2\to\C^2\), if
\[
\forall z\in\C^2,\qquad J_P(z)\ne0,
\]
then \(P\) is biholomorphic. Explicitly, there exists a holomorphic map \(Q:\C^2\to\C^2\) such that
\[
Q\circ P=\operatorname{id}_{\C^2}
\quad\text{and}\quad
P\circ Q=\operatorname{id}_{\C^2}.
\]
The supplied registry labels this statement open; its present literature status has not been independently verified.
\subsection{Short English statement}
Every polynomial map from \(\C^2\) to itself whose complex Jacobian determinant is nonzero everywhere is conjectured to be a bijection with a globally defined holomorphic inverse.
\subsection{Sources}
\begin{itemize}
\item Local UnsolvedMath registry, permanent ID \(800255\), source code \texttt{catalogue entry}, titled \emph{Several Complex Variables Problem List --- Problem 54}. The supplied record attributes the problem to Buzzard, Fornaess, and Krantz and gives the year 1995.
\item Buzzard, Fornaess, and Krantz, \emph{The Several Complex Variables Problem List} (1995), Problem 54. Supplied source URL: \url{https://arxiv.org/abs/math/9509202}. The record reports access on 2026-10-16 and extraction by \texttt{arxiv\_latex}; these are supplied provenance details, not a claim of independent source access or verification.
\end{itemize}
\end{document}
